% Source: https://artofproblemsolving.com/community/c3h1371399_trig_question
% Stable identity: aops-topic-1371399-q01
On the unit circle, the angle $90^\circ$ corresponds to the point $(0,1)$, so $\sin C=\sin90^\circ=\boxed{1}$. (The right-triangle definition ``opposite over hypotenuse'' only applies to the acute angles $A$ and $B$; for the right angle the unit-circle definition gives $1$, which also matches the law of sines: the side opposite $C$ is the hypotenuse $c$, and $\dfrac{c}{\sin C}=\dfrac{a}{\sin A}=c$.)

\begin{asy}
usepackage("amsmath");
size(3.5cm);
defaultpen(fontsize(10pt));
path ra(pair P,pair u,pair v,real s=0.3){return P+s*unit(u)--P+s*unit(u)+s*unit(v)--P+s*unit(v);}
pair Cc=(0,0),A=(4,0),B=(0,3);
draw(A--B--Cc--cycle,linewidth(0.9)); draw(ra(Cc,A,B,0.35));
dot(A);dot(B);dot(Cc);
label("$C$",Cc,SW); label("$A$",A,SE); label("$B$",B,N);
label("$a$",(B+Cc)/2,W); label("$b$",(A+Cc)/2,S); label("$c$",(A+B)/2,NE);
label("\mbox{}",(min(currentpicture,user=true)+max(currentpicture,user=true))/2,black); // keeps the text after the figure black when LaTeX typesets the labels inline
\end{asy}

With sides $a=BC$, $b=CA$ and hypotenuse $c=AB$, we have $\sin A=\dfrac {a}{c}$ and $\sin B=\dfrac {b}{c}$, so by the Pythagorean theorem
\[
\sin^2A+\sin^2B=\dfrac{a^2+b^2}{c^2}=1=\sin^2C .
\]
